\documentclass[a4paper]{article}

\usepackage{graphicx}
\usepackage{amsmath,amssymb}
\usepackage{minted}
\usepackage{multirow}
\usepackage{float}
\usepackage{todonotes}
\usepackage{forest}
\usepackage{xstring}
\usepackage{xcolor}
\usepackage[most]{tcolorbox}
\usepackage{xparse}
\usepackage{natbib}

\usetikzlibrary{graphs,graphdrawing}
\usegdlibrary{layered}

% for external graphviz
\input{/home/donkarlo/Dropbox/repo/nd_language_ai_project/src/nd_language_ai/formal/computer/programming/tex/latex/code_clips/external_graphviz.tex}

% for tags
\input{/home/donkarlo/Dropbox/repo/nd_language_ai_project/src/nd_language_ai/formal/computer/programming/tex/latex/parts/control_sequence/kind/semtag/semtag.tex}

% for links
\input{/home/donkarlo/Dropbox/repo/nd_language_ai_project/src/nd_language_ai/formal/computer/programming/tex/latex/code_clips/seehere.tex}

\title{Lokaladverben}
\date{}

\begin{document}
    \maketitle
    
    
    \section*{Adverbien wie \glqq dort\grqq}
        
        \textbf{dort} gehört grammatisch zu den \textbf{Adverbien}, genauer zu den \textbf{Lokaladverbien}.
        
        \subsection*{Was sind Lokaladverbien?}
            Lokaladverbien sind Adverbien des \textbf{Ortes}. Sie sagen:
            \begin{itemize}
                \item \textbf{wo} etwas ist,
                \item \textbf{wohin} etwas geht,
                \item \textbf{woher} etwas kommt.
            \end{itemize}
        
        \subsection*{Beispiele für Lokaladverbien}
            \begin{itemize}
                \item dort
                \item da
                \item hier
                \item draußen
                \item drinnen
                \item oben
                \item unten
                \item links
                \item rechts
                \item überall
                \item nirgends
            \end{itemize}
        
        \subsection*{Andere Gruppen von Adverbien}
            \begin{itemize}
                \item \textbf{Temporaladverbien} (Zeit): heute, gestern, bald, dann
                \item \textbf{Modaladverbien} (Art und Weise): gern, so, anders, absichtlich
                \item \textbf{Kausaladverbien} (Grund/Folge): deshalb, darum, daher, deswegen
            \end{itemize}
        
        \subsection*{Zusammenfassung}
            \begin{itemize}
                \item \textbf{dort} $\rightarrow$ Adverb
                \item genauer: \textbf{Lokaladverb}
            \end{itemize}
    %\bibliographystyle{plainnat}
    %\bibliography{/home/donkarlo/Dropbox/repo/nd_shared_research/bibliography/bibliography.bib}
\end{document}